\newpage
\phantomsection
\label{prf:metric-distance-to-set-zero-of-mem}
\LRAProofFor{thm:metric-distance-to-set-zero-of-mem}

\begin{remark*}[Return]
\hyperref[thm:metric-distance-to-set-zero-of-mem]{Return to Theorem}
\end{remark*}

\begin{theorem*}[Point-to-Set Distance is Zero at Points of the Set]
Let \((X,d)\) be a metric space, let \(A\subseteq X\), and let \(x\in A\).
Then \(d(x,A)=0\).
\end{theorem*}

\begin{proof}[Professional Standard Proof]
\LRAProofBodyStart
Since every element of \(D(x,A)\) is nonnegative, \(0\) is a lower bound for
\(D(x,A)\).  Hence \(0\leq d(x,A)\).  On the other hand, \(x\in A\), so
\[
  d(x,A)\leq d(x,x)=0.
\]
Thus \(d(x,A)\leq 0\) and \(0\leq d(x,A)\), so \(d(x,A)=0\).
\end{proof}

\begin{proof}[Detailed Learning Proof]
\LRAProofBodyStart
We prove the equality by proving both inequalities.

First, the point-to-set distance set \(D(x,A)\) consists of metric distances.
Every metric distance is nonnegative, so \(0\) is a lower bound for
\(D(x,A)\).  Since \(d(x,A)=\inf D(x,A)\), the infimum cannot be smaller than
this lower bound in the present nonnegative situation.  Hence
\[
  0\leq d(x,A).
\]

Second, because \(x\in A\), the witness-distance theorem gives
\[
  d(x,A)\leq d(x,x).
\]
The metric identity axiom gives \(d(x,x)=0\), so
\[
  d(x,A)\leq 0.
\]
Combining \(0\leq d(x,A)\) and \(d(x,A)\leq 0\), we conclude
\[
  d(x,A)=0.
\]
\end{proof}

\begin{remark*}[Proof structure]
Nonnegativity gives the lower bound \(0\leq d(x,A)\).  The witness \(x\in A\)
gives the upper bound \(d(x,A)\leq d(x,x)=0\).
\end{remark*}

\begin{dependencies}
\begin{itemize}
  \item \hyperref[def:metric-point-set-distance]{Definition: Distance from a Point to a Set}
  \item \hyperref[thm:metric-distance-set-bounded-below]{Theorem: Point-to-Set Distance Sets are Bounded Below}
  \item \hyperref[thm:metric-distance-to-set-le-distance-to-point-of-mem]{Theorem: Point-to-Set Distance is Bounded by Witness Distances}
\end{itemize}
\end{dependencies}

\clearpage
